\section{Conclusiones}
\label{vii:sec:conclusiones}

La amplitud torsional queda determinada por la acción espinorial y el
estado material, con una fórmula explícita y un sector en el que su
positividad se conserva. Este resultado permite obtener una contribución
energética proporcional a \(-a^{-6}\), un rebote homogéneo único bajo las
condiciones demostradas y, para polvo plano con
\(\Lambda_{\mathrm{eff}}=0\), una evolución regular para todo tiempo con
geodésicas causales de Levi--Civita completas. La secuencia enlaza un
observable cuártico del estado con una escala cosmológica mínima. Su
origen permanece en la construcción APP--TRIT--TPK: orientación, residuo,
acarreo y profundidad se transportan hasta las acciones, corrientes y
observables de la realización geométrica.

\subsection{Determinación de la amplitud y evolución cosmológica}

El resultado central de la realización espinorial es la determinación
de la amplitud torsional como funcional del estado material:
\[
 s_0[\varrho,V_c]
 =\frac{3\kappa c^2\hbar^2}{16}
   \frac{\gamma^2}{1+\gamma^2}
   \frac{\operatorname{Tr}
    (\varrho\widehat{\mathscr Q}_{\rm sp})}{V_c^2}.
\]
La fórmula procede sucesivamente de la representación interna de
Clifford, la corriente de espín, la eliminación de la contorsión,
el observable cuártico normalmente ordenado y las variaciones de
lapso y volumen. El factor cuadrático del volumen expresa la
dependencia de la densidad respecto del segundo momento de la
corriente. Una media de espín nula conserva un segundo momento
que puede contribuir a la geometría.

En el sector de ocupación tres se obtiene
\(\widehat{\mathscr Q}_{\rm sp}=4I\). Para estados de traza uno
soportados en ese sector, toda evolución unitaria que lo preserve
conserva su esperanza y determina una amplitud positiva.
El estado material y la celda comóvil individualizan la realización;
\(G\), \(\hbar\) y \(\gamma\) conservan sus antecedentes
estructurales. La constante reducida de Planck interviene cuadráticamente en
la interacción, Barbero--Immirzi determina el coeficiente del
operador de Holst y la gravitación convierte esa interacción en
respuesta geométrica. Esta composición explica conjuntamente
el valor funcional de la amplitud y la función de cada factor.

La ley de memoria \(M_n=\varphi^{-2n}\) y la lectura espacial
\(a_n/a_{\rm ref}=\varphi^{n/3}\) producen la dilución
\(a^{-6}\). En la realización espinorial de segundo momento
conservado, esa misma potencia aparece por la variación material.
Con \(\rho_0,s_0>0\), \(w<1\), curvatura espacial nula y
\(\Lambda_{\rm eff}=0\), el umbral de rebote es único y
transversal. Dos cotas diferenciadas preservan localmente su
existencia, unicidad y orientación temporal bajo perturbaciones.

El sector de polvo admite además una solución exacta para todo
tiempo. El factor de escala permanece separado de cero, la
curvatura es finita y las geodésicas causales de Levi--Civita
son completas en ambas direcciones. El resultado global pertenece
a esa realización explícita; la persistencia del rebote se
establece con las condiciones locales de su teorema. Ambas
pruebas muestran cómo la amplitud obtenida del estado se traduce
en una escala mínima y en una evolución de contracción y expansión.

\subsection{De la memoria transportada a la respuesta gravitatoria}

La normalización radial se obtiene después de la composición longitudinal
\(L_\alpha=(5759/23040)\alpha^{16}L_*\). El transporte normalizado
del archivo y la métrica positiva producen
\(R_X=L_\alpha\mathsf U X\mathsf U^*\); la energía conserva el mismo
carácter y su longitud conjugada satisface
\(H_X\Lambda_X=\hbar_{\rm ret}cI\). En la lectura de radio reducido,
estas igualdades determinan \(G=c^3L_\alpha^2/\hbar_{\rm ret}\), común
a los modos positivos de la realización. La prueba conserva sus reglas
longitudinal y métrica y la identificación física declarada. La reducción
conjunta de frontera y memoria mantiene el coeficiente; en la reducción
dinámica se transportan también el parámetro espectral y la norma temporal.
La expresión circular se recupera con
\(t_0=\pi_{\rm HMT}L_\alpha/(54c)\).

El espacio pentacomponente se obtiene mediante una proyección explícita
de los doce vértices y se transporta isométricamente a tensores
simétricos sin traza. La reducción transversal conserva dos componentes.
La secuencia \(12\to5\to5\to2\) distingue la incidencia inicial,
el portador gravitatorio, su representación tensorial y la lectura
transversal. El coeficiente \(G_a\) fija el acoplamiento común;
las orientaciones pertenecen al operador y a sus componentes.

En la misma sección de retorno, el coeficiente de la amplitud espinorial
admite la expresión equivalente
\[
 s_0[\varrho,V_c]
 =\frac{3\pi c\hbar_{\rm ret}L_\alpha^2}{2}
   \frac{\gamma^2}{1+\gamma^2}
   \frac{\operatorname{Tr}(\varrho\widehat{\mathscr Q}_{\rm sp})}{V_c^2}.
\]
Resulta de sustituir \(\kappa=8\pi L_\alpha^2/(\hbar_{\rm ret}c)\)
en el funcional ya demostrado. La longitud y la acción fijan así su
coeficiente, mientras el estado y el volumen conservan sus funciones
materiales en la solución cosmológica.

La pantalla de memoria proporciona una soldadura y una respuesta
energética compatibles con refinamiento. Su extensión minimizante,
construida a partir de la incidencia transportada, determina una
energía y un diferencial respecto de la conexión. El cálculo incluye
las variaciones de incidencia, frontera, peso y normalización. Al
pasar a la realización de Cartan, la cotetrada y el emparejamiento
material convierten ese diferencial en una corriente con dominio
geométrico explícito.

Las ecuaciones de conexión admiten una inversión algebraica de los
operadores de Holst y de torsión. En el sector material afín respecto
de la contorsión, esta inversión produce la corrección efectiva de
la acción. Para la extensión mínima, cuya corriente depende en
general de la propia conexión, se demuestra la ecuación estacionaria
completa, el criterio local de continuación mediante su Hessiano
y la fórmula integral de la acción efectiva. La acción espinorial y
la acción de extensión minimizante conservan, por tanto, sus argumentos
y sus leyes de variación propios. La primera proporciona el funcional de
amplitud evaluado en la solución cosmológica; la segunda describe la
respuesta variacional completa de la pantalla transportada.

\subsection{Información, horizontes y respuesta relacional}

Los balances residuales reúnen en una misma formulación la traza,
la parte sin traza y sus divergencias. Con acoplamientos constantes
en la coordenatización y fuente material conservada, el tensor residual
conserva covariantemente su balance total; la variación de su
traza determina el intercambio entre ambos sectores. La traza
constante caracteriza la conservación separada.

La información de horizonte utiliza la unidad areal,
la acción y la conversión térmica previamente construidas.
El estado reducido y su purificación distinguen información
accesible y correlaciones conservadas. En el balance esférico
desarrollado se verifican conjuntamente \(TS=E\),
\(T\,dS=2\,dE\) y \(S\,dT=-dE\); la dependencia de
la temperatura respecto del radio forma parte de esas identidades.
La constante de Boltzmann convierte el contenido informacional
en la sección térmica, mientras la relación acción--gravitación
determina la unidad de área de su lectura geométrica.

\Needspace{24\baselineskip}
La composición areal y cronológica determina
\(\mathcal T_{I/A}q_{\mathcal A}/\log3=h/(108t_0\log3)\).
En \(R=L_\alpha\), la temperatura cosmológica y la cronológica
están relacionadas por el factor \(\log3/(2\pi)\).
Para el Hamiltoniano de respuesta
\(\mathsf H_R=\tau_R\Delta_0\mathsf D_\partial\), los pesos
90 y 120 corresponden a distintos saltos energéticos con una escala
común. El balance finito es exactamente
\(\mathcal E_R-T_{\rm cos}\mathcal S_R
=\tau_R\mathcal D(\sigma\Vert\rho_0)\), y su término tangente
es la primera ley a radio fijo. La sección relativa construida desde
ese estado determina un único transductor positivo; la identificación
de un lector adicional se reduce a comprobar su valor sobre la sección
cronológica de~\eqref{viii:eq:criterio-una-seccion}. Las isometrías que
conservan el registro transportan este balance y sus observables.

El observador se define mediante operaciones de orientación,
trivialización y restricción de observables. La respuesta
autoinercial conserva su factorización por la información global
y hace explícito el defecto producido por la proyección.
Finalmente, en el canal cronológico realizado, la diferencia covariante
construida con el transporte de rutas admite, para fuentes de soporte
finito, operadores de Green retardados y avanzados. La compatibilidad de ambas lecturas
en la frontera delimita un espacio de fuentes y una proyección
que interviene en la respuesta variacional. La bidireccionalidad
queda expresada por operadores, momentos y condiciones de frontera
que pueden comprobarse dentro del artículo.

\enlargethispage{2\baselineskip}
La relación decisiva es que la información conservada por el estado
interviene en la respuesta geométrica mediante una corriente y su segundo
momento. La eliminación de la contorsión y la variación métrica convierten
ese observable en una densidad y una presión; la solución de polvo muestra
cómo ambas determinan una escala mínima y una evolución causal completa.
Los balances y la lectura del observador conservan, a su vez, la
distinción entre información accesible y correlación global. En los
dominios construidos, memoria, fuente material y geometría forman así una
cadena explícita de determinaciones, desde el núcleo aritmético hasta la
evolución gravitatoria y su contenido informacional.
